% Options for packages loaded elsewhere
% Options for packages loaded elsewhere
\PassOptionsToPackage{unicode}{hyperref}
\PassOptionsToPackage{hyphens}{url}
\PassOptionsToPackage{dvipsnames,svgnames,x11names}{xcolor}
%
\documentclass[
  letterpaper,
  DIV=11,
  numbers=noendperiod]{scrartcl}
\usepackage{xcolor}
\usepackage{amsmath,amssymb}
\setcounter{secnumdepth}{5}
\usepackage{iftex}
\ifPDFTeX
  \usepackage[T1]{fontenc}
  \usepackage[utf8]{inputenc}
  \usepackage{textcomp} % provide euro and other symbols
\else % if luatex or xetex
  \usepackage{unicode-math} % this also loads fontspec
  \defaultfontfeatures{Scale=MatchLowercase}
  \defaultfontfeatures[\rmfamily]{Ligatures=TeX,Scale=1}
\fi
\usepackage{lmodern}
\ifPDFTeX\else
  % xetex/luatex font selection
\fi
% Use upquote if available, for straight quotes in verbatim environments
\IfFileExists{upquote.sty}{\usepackage{upquote}}{}
\IfFileExists{microtype.sty}{% use microtype if available
  \usepackage[]{microtype}
  \UseMicrotypeSet[protrusion]{basicmath} % disable protrusion for tt fonts
}{}
\makeatletter
\@ifundefined{KOMAClassName}{% if non-KOMA class
  \IfFileExists{parskip.sty}{%
    \usepackage{parskip}
  }{% else
    \setlength{\parindent}{0pt}
    \setlength{\parskip}{6pt plus 2pt minus 1pt}}
}{% if KOMA class
  \KOMAoptions{parskip=half}}
\makeatother
% Make \paragraph and \subparagraph free-standing
\makeatletter
\ifx\paragraph\undefined\else
  \let\oldparagraph\paragraph
  \renewcommand{\paragraph}{
    \@ifstar
      \xxxParagraphStar
      \xxxParagraphNoStar
  }
  \newcommand{\xxxParagraphStar}[1]{\oldparagraph*{#1}\mbox{}}
  \newcommand{\xxxParagraphNoStar}[1]{\oldparagraph{#1}\mbox{}}
\fi
\ifx\subparagraph\undefined\else
  \let\oldsubparagraph\subparagraph
  \renewcommand{\subparagraph}{
    \@ifstar
      \xxxSubParagraphStar
      \xxxSubParagraphNoStar
  }
  \newcommand{\xxxSubParagraphStar}[1]{\oldsubparagraph*{#1}\mbox{}}
  \newcommand{\xxxSubParagraphNoStar}[1]{\oldsubparagraph{#1}\mbox{}}
\fi
\makeatother

\usepackage{color}
\usepackage{fancyvrb}
\newcommand{\VerbBar}{|}
\newcommand{\VERB}{\Verb[commandchars=\\\{\}]}
\DefineVerbatimEnvironment{Highlighting}{Verbatim}{commandchars=\\\{\}}
% Add ',fontsize=\small' for more characters per line
\usepackage{framed}
\definecolor{shadecolor}{RGB}{241,243,245}
\newenvironment{Shaded}{\begin{snugshade}}{\end{snugshade}}
\newcommand{\AlertTok}[1]{\textcolor[rgb]{0.68,0.00,0.00}{#1}}
\newcommand{\AnnotationTok}[1]{\textcolor[rgb]{0.37,0.37,0.37}{#1}}
\newcommand{\AttributeTok}[1]{\textcolor[rgb]{0.40,0.45,0.13}{#1}}
\newcommand{\BaseNTok}[1]{\textcolor[rgb]{0.68,0.00,0.00}{#1}}
\newcommand{\BuiltInTok}[1]{\textcolor[rgb]{0.00,0.23,0.31}{#1}}
\newcommand{\CharTok}[1]{\textcolor[rgb]{0.13,0.47,0.30}{#1}}
\newcommand{\CommentTok}[1]{\textcolor[rgb]{0.37,0.37,0.37}{#1}}
\newcommand{\CommentVarTok}[1]{\textcolor[rgb]{0.37,0.37,0.37}{\textit{#1}}}
\newcommand{\ConstantTok}[1]{\textcolor[rgb]{0.56,0.35,0.01}{#1}}
\newcommand{\ControlFlowTok}[1]{\textcolor[rgb]{0.00,0.23,0.31}{\textbf{#1}}}
\newcommand{\DataTypeTok}[1]{\textcolor[rgb]{0.68,0.00,0.00}{#1}}
\newcommand{\DecValTok}[1]{\textcolor[rgb]{0.68,0.00,0.00}{#1}}
\newcommand{\DocumentationTok}[1]{\textcolor[rgb]{0.37,0.37,0.37}{\textit{#1}}}
\newcommand{\ErrorTok}[1]{\textcolor[rgb]{0.68,0.00,0.00}{#1}}
\newcommand{\ExtensionTok}[1]{\textcolor[rgb]{0.00,0.23,0.31}{#1}}
\newcommand{\FloatTok}[1]{\textcolor[rgb]{0.68,0.00,0.00}{#1}}
\newcommand{\FunctionTok}[1]{\textcolor[rgb]{0.28,0.35,0.67}{#1}}
\newcommand{\ImportTok}[1]{\textcolor[rgb]{0.00,0.46,0.62}{#1}}
\newcommand{\InformationTok}[1]{\textcolor[rgb]{0.37,0.37,0.37}{#1}}
\newcommand{\KeywordTok}[1]{\textcolor[rgb]{0.00,0.23,0.31}{\textbf{#1}}}
\newcommand{\NormalTok}[1]{\textcolor[rgb]{0.00,0.23,0.31}{#1}}
\newcommand{\OperatorTok}[1]{\textcolor[rgb]{0.37,0.37,0.37}{#1}}
\newcommand{\OtherTok}[1]{\textcolor[rgb]{0.00,0.23,0.31}{#1}}
\newcommand{\PreprocessorTok}[1]{\textcolor[rgb]{0.68,0.00,0.00}{#1}}
\newcommand{\RegionMarkerTok}[1]{\textcolor[rgb]{0.00,0.23,0.31}{#1}}
\newcommand{\SpecialCharTok}[1]{\textcolor[rgb]{0.37,0.37,0.37}{#1}}
\newcommand{\SpecialStringTok}[1]{\textcolor[rgb]{0.13,0.47,0.30}{#1}}
\newcommand{\StringTok}[1]{\textcolor[rgb]{0.13,0.47,0.30}{#1}}
\newcommand{\VariableTok}[1]{\textcolor[rgb]{0.07,0.07,0.07}{#1}}
\newcommand{\VerbatimStringTok}[1]{\textcolor[rgb]{0.13,0.47,0.30}{#1}}
\newcommand{\WarningTok}[1]{\textcolor[rgb]{0.37,0.37,0.37}{\textit{#1}}}

\usepackage{longtable,booktabs,array}
\newcounter{none} % for unnumbered tables
\usepackage{calc} % for calculating minipage widths
% Correct order of tables after \paragraph or \subparagraph
\usepackage{etoolbox}
\makeatletter
\patchcmd\longtable{\par}{\if@noskipsec\mbox{}\fi\par}{}{}
\makeatother
% Allow footnotes in longtable head/foot
\IfFileExists{footnotehyper.sty}{\usepackage{footnotehyper}}{\usepackage{footnote}}
\makesavenoteenv{longtable}
\usepackage{graphicx}
\makeatletter
\newsavebox\pandoc@box
\newcommand*\pandocbounded[1]{% scales image to fit in text height/width
  \sbox\pandoc@box{#1}%
  \Gscale@div\@tempa{\textheight}{\dimexpr\ht\pandoc@box+\dp\pandoc@box\relax}%
  \Gscale@div\@tempb{\linewidth}{\wd\pandoc@box}%
  \ifdim\@tempb\p@<\@tempa\p@\let\@tempa\@tempb\fi% select the smaller of both
  \ifdim\@tempa\p@<\p@\scalebox{\@tempa}{\usebox\pandoc@box}%
  \else\usebox{\pandoc@box}%
  \fi%
}
% Set default figure placement to htbp
\def\fps@figure{htbp}
\makeatother





\setlength{\emergencystretch}{3em} % prevent overfull lines

\providecommand{\tightlist}{%
  \setlength{\itemsep}{0pt}\setlength{\parskip}{0pt}}



 


% load packages
\usepackage{geometry}
\usepackage{xcolor}
\usepackage{eso-pic}
\usepackage{fancyhdr}
\usepackage{sectsty}
\usepackage{fontspec}
\usepackage{titlesec}

%% Set page size with a wider right margin
\geometry{a4paper, total={170mm,257mm}, left=20mm, top=20mm, bottom=20mm, right=50mm}

%% Let's define some colours
\definecolor{light}{HTML}{E6E6FA}
\definecolor{highlight}{HTML}{800080}
\definecolor{dark}{HTML}{330033}

%% Let's add the border on the right hand side 
\AddToShipoutPicture{% 
    \AtPageLowerLeft{% 
        \put(\LenToUnit{\dimexpr\paperwidth-3cm},0){% 
            \color{light}\rule{3cm}{\LenToUnit\paperheight}%
          }%
     }%
     % logo
    \AtPageLowerLeft{% start the bar at the bottom right of the page
        \put(\LenToUnit{\dimexpr\paperwidth-2.25cm},27.2cm){% move it to the top right
            \color{light}\includegraphics[width=1.5cm]{_extensions/nrennie/PrettyPDF/logo.png}
          }%
     }%
}

%% Style the page number
\fancypagestyle{mystyle}{
  \fancyhf{}
  \renewcommand\headrulewidth{0pt}
  \fancyfoot[R]{\thepage}
  \fancyfootoffset{3.5cm}
}
\setlength{\footskip}{20pt}

%% style the chapter/section fonts
\chapterfont{\color{dark}\fontsize{20}{16.8}\selectfont}
\sectionfont{\color{dark}\fontsize{20}{16.8}\selectfont}
\subsectionfont{\color{dark}\fontsize{14}{16.8}\selectfont}
\titleformat{\subsection}
  {\sffamily\Large\bfseries}{\thesection}{1em}{}[{\titlerule[0.8pt]}]
  
% left align title
\makeatletter
\renewcommand{\maketitle}{\bgroup\setlength{\parindent}{0pt}
\begin{flushleft}
  {\sffamily\huge\textbf{\MakeUppercase{\@title}}} \vspace{0.3cm} \newline
  {\Large {\@subtitle}} \newline
  \@author
\end{flushleft}\egroup
}
\makeatother

%% Use some custom fonts
\setsansfont{Ubuntu}[
    Path=_extensions/nrennie/PrettyPDF/Ubuntu/,
    Scale=0.9,
    Extension = .ttf,
    UprightFont=*-Regular,
    BoldFont=*-Bold,
    ItalicFont=*-Italic,
    ]

\setmainfont{Ubuntu}[
    Path=_extensions/nrennie/PrettyPDF/Ubuntu/,
    Scale=0.9,
    Extension = .ttf,
    UprightFont=*-Regular,
    BoldFont=*-Bold,
    ItalicFont=*-Italic,
    ]
\KOMAoption{captions}{tableheading}
\makeatletter
\@ifpackageloaded{tcolorbox}{}{\usepackage[skins,breakable]{tcolorbox}}
\@ifpackageloaded{fontawesome5}{}{\usepackage{fontawesome5}}
\definecolor{quarto-callout-color}{HTML}{909090}
\definecolor{quarto-callout-note-color}{HTML}{0758E5}
\definecolor{quarto-callout-important-color}{HTML}{CC1914}
\definecolor{quarto-callout-warning-color}{HTML}{EB9113}
\definecolor{quarto-callout-tip-color}{HTML}{00A047}
\definecolor{quarto-callout-caution-color}{HTML}{FC5300}
\definecolor{quarto-callout-color-frame}{HTML}{acacac}
\definecolor{quarto-callout-note-color-frame}{HTML}{4582ec}
\definecolor{quarto-callout-important-color-frame}{HTML}{d9534f}
\definecolor{quarto-callout-warning-color-frame}{HTML}{f0ad4e}
\definecolor{quarto-callout-tip-color-frame}{HTML}{02b875}
\definecolor{quarto-callout-caution-color-frame}{HTML}{fd7e14}
\makeatother
\makeatletter
\@ifpackageloaded{caption}{}{\usepackage{caption}}
\AtBeginDocument{%
\ifdefined\contentsname
  \renewcommand*\contentsname{Table of contents}
\else
  \newcommand\contentsname{Table of contents}
\fi
\ifdefined\listfigurename
  \renewcommand*\listfigurename{List of Figures}
\else
  \newcommand\listfigurename{List of Figures}
\fi
\ifdefined\listtablename
  \renewcommand*\listtablename{List of Tables}
\else
  \newcommand\listtablename{List of Tables}
\fi
\ifdefined\figurename
  \renewcommand*\figurename{Figure}
\else
  \newcommand\figurename{Figure}
\fi
\ifdefined\tablename
  \renewcommand*\tablename{Table}
\else
  \newcommand\tablename{Table}
\fi
}
\@ifpackageloaded{float}{}{\usepackage{float}}
\floatstyle{ruled}
\@ifundefined{c@chapter}{\newfloat{codelisting}{h}{lop}}{\newfloat{codelisting}{h}{lop}[chapter]}
\floatname{codelisting}{Listing}
\newcommand*\listoflistings{\listof{codelisting}{List of Listings}}
\makeatother
\makeatletter
\makeatother
\makeatletter
\@ifpackageloaded{caption}{}{\usepackage{caption}}
\@ifpackageloaded{subcaption}{}{\usepackage{subcaption}}
\makeatother
\makeatletter
\@ifpackageloaded{tcolorbox}{}{\usepackage[skins,breakable]{tcolorbox}}
\makeatother
\makeatletter
\@ifundefined{shadecolor}{\definecolor{shadecolor}{rgb}{.97, .97, .97}}{}
\makeatother
\makeatletter
\@ifundefined{codebgcolor}{\definecolor{codebgcolor}{named}{light}}{}
\makeatother
\makeatletter
\ifdefined\Shaded\renewenvironment{Shaded}{\begin{tcolorbox}[boxrule=0pt, breakable, colback={codebgcolor}, enhanced, frame hidden, sharp corners]}{\end{tcolorbox}}\fi
\makeatother
\usepackage{bookmark}
\IfFileExists{xurl.sty}{\usepackage{xurl}}{} % add URL line breaks if available
\urlstyle{same}
\hypersetup{
  pdftitle={Practical - Geostatistics},
  colorlinks=true,
  linkcolor={highlight},
  filecolor={Maroon},
  citecolor={Blue},
  urlcolor={highlight},
  pdfcreator={LaTeX via pandoc}}


\title{Practical - Geostatistics}
\author{}
\date{}
\begin{document}
\maketitle

\pagestyle{mystyle}


\section{Geostatistics}\label{geostatistics}

In this practical we are going to fit a geostatistical model. We will:

\begin{itemize}
\item
  Learn how to fit a geostatistical model in \texttt{inlabru}
\item
  Learn how to build a mesh for the SPDE representation
\item
  Learn how to add spatial covariates to the model
\item
  Learn how to do predictions
\item
  Learn how to simulate from the fitted model
\end{itemize}

\begin{center}\rule{0.5\linewidth}{0.5pt}\end{center}

Libraries to load:

\begin{Shaded}
\begin{Highlighting}[]
\FunctionTok{library}\NormalTok{(dplyr)}
\FunctionTok{library}\NormalTok{(INLA)}
\FunctionTok{library}\NormalTok{(inlabru)}
\FunctionTok{library}\NormalTok{(fmesher)}
\FunctionTok{library}\NormalTok{(sf)}
\FunctionTok{library}\NormalTok{(terra)}


\CommentTok{\# load some libraries to generate nice map plots}
\FunctionTok{library}\NormalTok{(scico)}
\FunctionTok{library}\NormalTok{(ggplot2)}
\FunctionTok{library}\NormalTok{(patchwork)}
\FunctionTok{library}\NormalTok{(mapview)}
\FunctionTok{library}\NormalTok{(tidyterra)}
\end{Highlighting}
\end{Shaded}

\subsection{The data}\label{the-data}

We will explore data on the Pacific Cod (\emph{Gadus macrocephalus})
from a trawl survey in Queen Charlotte Sound. The \texttt{pcod} dataset
is available from the \texttt{sdmTMB} package or can be downloaded by
clicking the button below:

The dataset contains the presence/absence records of the Pacific Cod
during each surveys along with the biomass density of Pacific cod in the
area swept (kg/Km\(^2\)). The \texttt{qcs\_grid} data contain the depth
values stored as \(2\times 2\) km grid for Queen Charlotte Sound. The
data contains presence/absence data from 2003 to 2017. In this practical
we only consider year 2003.

We first load the dataset and select the year of interest

\begin{Shaded}
\begin{Highlighting}[]
\FunctionTok{load}\NormalTok{(here}\SpecialCharTok{::}\FunctionTok{here}\NormalTok{(}\StringTok{"datasets/pcod.RData"}\NormalTok{))}
\end{Highlighting}
\end{Shaded}

Then, we create an \texttt{sf} object and assign the rough coordinate
reference to it:

\begin{Shaded}
\begin{Highlighting}[]
\NormalTok{pcod\_sf }\OtherTok{=}   \FunctionTok{st\_as\_sf}\NormalTok{(pcod\_df, }\AttributeTok{coords =} \FunctionTok{c}\NormalTok{(}\StringTok{"lon"}\NormalTok{,}\StringTok{"lat"}\NormalTok{), }\AttributeTok{crs =} \DecValTok{4326}\NormalTok{)}
\NormalTok{pcod\_sf }\OtherTok{=} \FunctionTok{st\_transform}\NormalTok{(pcod\_sf,}
\AttributeTok{crs =} \StringTok{"+proj=utm +zone=9 +datum=WGS84 +no\_defs +type=crs +units=km"}\NormalTok{ )}
\end{Highlighting}
\end{Shaded}

We convert the covariate into a raster and assign the same coordinate
reference:

\begin{Shaded}
\begin{Highlighting}[]
\NormalTok{depth\_r }\OtherTok{\textless{}{-}} \FunctionTok{rast}\NormalTok{(qcs\_grid, }\AttributeTok{type =} \StringTok{"xyz"}\NormalTok{)}
\FunctionTok{crs}\NormalTok{(depth\_r) }\OtherTok{\textless{}{-}} \FunctionTok{crs}\NormalTok{(pcod\_sf)}
\end{Highlighting}
\end{Shaded}

Finally, we can plot our dataset. Note that to plot the raster we need
to also upload the \texttt{tidyterra} library.

\begin{Shaded}
\begin{Highlighting}[]
\FunctionTok{ggplot}\NormalTok{()}\SpecialCharTok{+} 
  \FunctionTok{geom\_spatraster}\NormalTok{(}\AttributeTok{data=}\NormalTok{depth\_r}\SpecialCharTok{$}\NormalTok{depth)}\SpecialCharTok{+}
  \FunctionTok{geom\_sf}\NormalTok{(}\AttributeTok{data=}\NormalTok{pcod\_sf,}\FunctionTok{aes}\NormalTok{(}\AttributeTok{color=}\FunctionTok{factor}\NormalTok{(present))) }\SpecialCharTok{+}
    \FunctionTok{scale\_color\_manual}\NormalTok{(}\AttributeTok{name=}\StringTok{"Occupancy status for the Pacific Cod"}\NormalTok{,}
                     \AttributeTok{values =} \FunctionTok{c}\NormalTok{(}\StringTok{"black"}\NormalTok{,}\StringTok{"orange"}\NormalTok{),}
                     \AttributeTok{labels=} \FunctionTok{c}\NormalTok{(}\StringTok{"Absence"}\NormalTok{,}\StringTok{"Presence"}\NormalTok{))}\SpecialCharTok{+}
  \FunctionTok{scale\_fill\_scico}\NormalTok{(}\AttributeTok{name =} \StringTok{"Depth"}\NormalTok{,}
                   \AttributeTok{palette =} \StringTok{"nuuk"}\NormalTok{,}
                   \AttributeTok{na.value =} \StringTok{"transparent"}\NormalTok{ ) }\SpecialCharTok{+} \FunctionTok{xlab}\NormalTok{(}\StringTok{""}\NormalTok{) }\SpecialCharTok{+} \FunctionTok{ylab}\NormalTok{(}\StringTok{""}\NormalTok{)}
\end{Highlighting}
\end{Shaded}

\pandocbounded{\includegraphics[keepaspectratio]{Geostat_ex_files/figure-pdf/unnamed-chunk-6-1.pdf}}

\subsection{The model}\label{the-model}

We first fit a simple model where we consider the observation as
Bernoulli and where the linear predictor contains only one intercept and
the GR field defined through the SPDE approach. The model is defined as:

\textbf{Stage 1} Model for the response

\[
y(s)|\eta(s)\sim\text{Binom}(1, p(s))
\] \textbf{Stage 2} Latent field model

\[
\eta(s) = \text{logit}(p(s)) = \beta_0 + \omega(s)
\]

with

\[
\omega(s)\sim \text{  GF with range } \rho\  \text{ and maginal variance }\ \sigma^2
\]

\textbf{Stage 3} Hyperparameters

The hyperparameters of the model are \(\rho\) and \(\sigma\)

\textbf{NOTE} In this case the linear predictor \(\eta\) consists of two
components!!

\subsection{The workflow}\label{the-workflow}

When fitting a geostatistical model we need to go through the following
steps:

\begin{enumerate}
\def\labelenumi{\arabic{enumi}.}
\item
  Build the mesh
\item
  Define the SPDE representation of the spatial GF. This includes
  defining the priors for the range and sd of the spatial GF
\item
  Define the \emph{components} of the linear predictor. This includes
  the spatial GF and all potential covariates
\item
  Define the observation model using the \texttt{bru\_obs()} function
\item
  Run the model using the \texttt{bru()} function
\end{enumerate}

\subsubsection{1. Building the mesh}\label{building-the-mesh}

The first task, when dealing with geostatistical models in
\texttt{inlabru} is to build the mesh that covers the area of interest.
For this purpose we use the function \texttt{fm\_messh\_2d}.

One way to build the mesh is to start from the locations where we have
observations, these are contained in the dataset \texttt{pcod\_sf}.

\begin{Shaded}
\begin{Highlighting}[]
\NormalTok{mesh }\OtherTok{=} \FunctionTok{fm\_mesh\_2d}\NormalTok{(}\AttributeTok{loc =}\NormalTok{ pcod\_sf,           }\CommentTok{\# Build the mesh}
                  \AttributeTok{max.edge =} \FunctionTok{c}\NormalTok{(}\DecValTok{10}\NormalTok{,}\DecValTok{20}\NormalTok{),     }\CommentTok{\# The largest allowed triangle edge length.}
                  \AttributeTok{offset =} \FunctionTok{c}\NormalTok{(}\DecValTok{5}\NormalTok{,}\DecValTok{50}\NormalTok{))       }\CommentTok{\# The automatic extension distance}
\FunctionTok{ggplot}\NormalTok{() }\SpecialCharTok{+} \FunctionTok{gg}\NormalTok{(mesh) }\SpecialCharTok{+} \FunctionTok{geom\_sf}\NormalTok{(}\AttributeTok{data=}\NormalTok{ pcod\_sf, }\FunctionTok{aes}\NormalTok{(}\AttributeTok{color =} \FunctionTok{factor}\NormalTok{(present)), }\AttributeTok{size =} \FloatTok{0.1}\NormalTok{) }\SpecialCharTok{+} \FunctionTok{xlab}\NormalTok{(}\StringTok{""}\NormalTok{) }\SpecialCharTok{+} \FunctionTok{ylab}\NormalTok{(}\StringTok{""}\NormalTok{)}
\end{Highlighting}
\end{Shaded}

\pandocbounded{\includegraphics[keepaspectratio]{Geostat_ex_files/figure-pdf/unnamed-chunk-7-1.pdf}}

As you can see from the plot above, some of the locations are very close
to each other, this causes some very small triangles. This can be
avoided using the option \texttt{cutoff\ =} which collapses the
locations that are closer than a cutoff (those points are collapsed in
the mesh construction but, of course, not when it come to estimaation.)

\begin{Shaded}
\begin{Highlighting}[]
\NormalTok{mesh }\OtherTok{=} \FunctionTok{fm\_mesh\_2d}\NormalTok{(}\AttributeTok{loc =}\NormalTok{ pcod\_sf,           }\CommentTok{\# Build the mesh}
                  \AttributeTok{cutoff =} \DecValTok{2}\NormalTok{,}
                  \AttributeTok{max.edge =} \FunctionTok{c}\NormalTok{(}\DecValTok{10}\NormalTok{,}\DecValTok{20}\NormalTok{),     }\CommentTok{\# The largest allowed triangle edge length.}
                  \AttributeTok{offset =} \FunctionTok{c}\NormalTok{(}\DecValTok{5}\NormalTok{,}\DecValTok{50}\NormalTok{))       }\CommentTok{\# The automatic extension distance}
\FunctionTok{ggplot}\NormalTok{() }\SpecialCharTok{+} \FunctionTok{gg}\NormalTok{(mesh) }\SpecialCharTok{+} \FunctionTok{geom\_sf}\NormalTok{(}\AttributeTok{data=}\NormalTok{ pcod\_sf, }\FunctionTok{aes}\NormalTok{(}\AttributeTok{color =} \FunctionTok{factor}\NormalTok{(present)), }\AttributeTok{size =} \FloatTok{0.1}\NormalTok{) }\SpecialCharTok{+} \FunctionTok{xlab}\NormalTok{(}\StringTok{""}\NormalTok{) }\SpecialCharTok{+} \FunctionTok{ylab}\NormalTok{(}\StringTok{""}\NormalTok{)}
\end{Highlighting}
\end{Shaded}

\pandocbounded{\includegraphics[keepaspectratio]{Geostat_ex_files/figure-pdf/unnamed-chunk-8-1.pdf}}

Key parameters in mesh construction include: \texttt{max.edge} for
maximum triangle edge lengths, \texttt{offset} for inner and outer
extensions (to prevent edge effects), and cutoff to avoid overly small
triangles in clustered areas.

\begin{tcolorbox}[enhanced jigsaw, arc=.35mm, bottomrule=.15mm, bottomtitle=1mm, breakable, colback=white, colbacktitle=quarto-callout-note-color!10!white, colframe=quarto-callout-note-color-frame, coltitle=black, left=2mm, leftrule=.75mm, opacityback=0, opacitybacktitle=0.6, rightrule=.15mm, title=\textcolor{quarto-callout-note-color}{\faInfo}\hspace{0.5em}{Note}, titlerule=0mm, toprule=.15mm, toptitle=1mm]

\textbf{General guidelines for creating the mesh}

\begin{enumerate}
\def\labelenumi{\arabic{enumi}.}
\tightlist
\item
  Create triangulation meshes with \texttt{fm\_mesh\_2d()}
\item
  Move undesired boundary effects away from the domain of interest by
  extending to a smooth external boundary
\item
  Use a coarser resolution in the extension to reduce computational cost
  (\texttt{max.edge=c(inner,\ outer)})
\item
  Use a fine resolution (subject to available computational resources)
  for the domain of interest (inner correlation range) and filter out
  small input point clusters (0 \textless{} \texttt{cutoff} \textless{}
  inner)
\item
  Coastlines and similar can be added to the domain specification in
  \texttt{fm\_mesh\_2d()} through the \texttt{boundary} argument.
\end{enumerate}

\end{tcolorbox}

\begin{tcolorbox}[enhanced jigsaw, arc=.35mm, bottomrule=.15mm, bottomtitle=1mm, breakable, colback=white, colbacktitle=quarto-callout-warning-color!10!white, colframe=quarto-callout-warning-color-frame, coltitle=black, left=2mm, leftrule=.75mm, opacityback=0, opacitybacktitle=0.6, rightrule=.15mm, title={Task}, titlerule=0mm, toprule=.15mm, toptitle=1mm]

Look at the documentation for the \texttt{fm\_mesh\_2d} function typing

\begin{Shaded}
\begin{Highlighting}[]
\NormalTok{?fm\_mesh\_2d}
\end{Highlighting}
\end{Shaded}

play around with the different options and create different meshes.

The \emph{rule of thumb} is that your mesh should be:

\begin{itemize}
\tightlist
\item
  fine enough to well represent the spatial variability of your process,
  but not too fine in order to avoid computation burden
\item
  the triangles should be regular, avoid long and thin triangles.
\item
  The mesh should contain a buffer around your area of interest (this is
  what is defined in the \texttt{offset} option) in order to avoid
  boundary artefact in the estimated variance.
\end{itemize}

\end{tcolorbox}

\subsubsection{2. Define the SPDE representation of the spatial
GF}\label{define-the-spde-representation-of-the-spatial-gf}

To define the SPDE representation of the spatial GF we use the function
\texttt{inla.spde2.pcmatern}. This takes as input the mesh we have
defined and the PC-priors definition for \(\rho\) and \(\sigma\) (the
range and the marginal standard deviation of the field).

PC priors Gaussian Random field are defined in (Fuglstad et al.~2018).
From a practical perspective for the range \(\rho\) you need to define
two paramters \(\rho_0\) and \(p_{\rho}\) such that you believe it is
reasonable that

\[
P(\rho<\rho_0)=p_{\rho}
\]

while for the margianal variance \(\sigma\) you need to define two
parameters \(\sigma_0\) and \(p_{\sigma}\) such that you believe it is
reasonable that

\[
P(\sigma<\sigma_0)=p_{\sigma}
\]

You can use the following function to compute and plot the prior
distributions for the range and sd of the Matern field.

\begin{Shaded}
\begin{Highlighting}[]
\NormalTok{dens\_prior\_range }\OtherTok{=} \ControlFlowTok{function}\NormalTok{(rho\_0, p\_alpha)}
\NormalTok{\{}
  \CommentTok{\# compute the density of the PC prior for the}
  \CommentTok{\# range rho of the Matern field}
  \CommentTok{\# rho\_0 and p\_alpha are defined such that}
  \CommentTok{\# P(rho\textless{}rho\_0) = p\_alpha}
\NormalTok{  rho }\OtherTok{=} \FunctionTok{seq}\NormalTok{(}\DecValTok{0}\NormalTok{, rho\_0}\SpecialCharTok{*}\DecValTok{10}\NormalTok{, }\AttributeTok{length.out =}\DecValTok{100}\NormalTok{)}
\NormalTok{  alpha1\_tilde }\OtherTok{=} \SpecialCharTok{{-}}\FunctionTok{log}\NormalTok{(p\_alpha) }\SpecialCharTok{*}\NormalTok{ rho\_0}
\NormalTok{  dens\_rho }\OtherTok{=}\NormalTok{  alpha1\_tilde }\SpecialCharTok{/}\NormalTok{ rho}\SpecialCharTok{\^{}}\DecValTok{2} \SpecialCharTok{*} \FunctionTok{exp}\NormalTok{(}\SpecialCharTok{{-}}\NormalTok{alpha1\_tilde }\SpecialCharTok{/}\NormalTok{ rho)}
  \FunctionTok{return}\NormalTok{(}\FunctionTok{data.frame}\NormalTok{(}\AttributeTok{x =}\NormalTok{ rho, }\AttributeTok{y =}\NormalTok{ dens\_rho))}
\NormalTok{\}}

\NormalTok{dens\_prior\_sd }\OtherTok{=} \ControlFlowTok{function}\NormalTok{(sigma\_0, p\_sigma)}
\NormalTok{\{}
  \CommentTok{\# compute the density of the PC prior for the}
  \CommentTok{\# sd sigma of the Matern field}
  \CommentTok{\# sigma\_0 and p\_sigma are defined such that}
  \CommentTok{\# P(sigma\textgreater{}sigma\_0) = p\_sigma}
\NormalTok{  sigma }\OtherTok{=} \FunctionTok{seq}\NormalTok{(}\DecValTok{0}\NormalTok{, sigma\_0}\SpecialCharTok{*}\DecValTok{10}\NormalTok{, }\AttributeTok{length.out =}\DecValTok{100}\NormalTok{)}
\NormalTok{  alpha2\_tilde }\OtherTok{=} \SpecialCharTok{{-}}\FunctionTok{log}\NormalTok{(p\_sigma)}\SpecialCharTok{/}\NormalTok{sigma\_0}
\NormalTok{  dens\_sigma }\OtherTok{=}\NormalTok{ alpha2\_tilde}\SpecialCharTok{*} \FunctionTok{exp}\NormalTok{(}\SpecialCharTok{{-}}\NormalTok{alpha2\_tilde }\SpecialCharTok{*}\NormalTok{ sigma) }
  \FunctionTok{return}\NormalTok{(}\FunctionTok{data.frame}\NormalTok{(}\AttributeTok{x =}\NormalTok{ sigma, }\AttributeTok{y =}\NormalTok{ dens\_sigma))}
\NormalTok{\}}
\end{Highlighting}
\end{Shaded}

Here are some alternatives for defining priors for our model

\begin{Shaded}
\begin{Highlighting}[]
\NormalTok{spde\_model1 }\OtherTok{=}  \FunctionTok{inla.spde2.pcmatern}\NormalTok{(mesh,}
                                  \AttributeTok{prior.sigma =} \FunctionTok{c}\NormalTok{(.}\DecValTok{1}\NormalTok{, }\FloatTok{0.5}\NormalTok{),}
                                  \AttributeTok{prior.range =} \FunctionTok{c}\NormalTok{(}\DecValTok{30}\NormalTok{, }\FloatTok{0.5}\NormalTok{))}
\NormalTok{spde\_model2 }\OtherTok{=}  \FunctionTok{inla.spde2.pcmatern}\NormalTok{(mesh,}
                                  \AttributeTok{prior.sigma =} \FunctionTok{c}\NormalTok{(}\DecValTok{10}\NormalTok{, }\FloatTok{0.5}\NormalTok{),}
                                  \AttributeTok{prior.range =} \FunctionTok{c}\NormalTok{(}\DecValTok{1000}\NormalTok{, }\FloatTok{0.5}\NormalTok{))}
\NormalTok{spde\_model3 }\OtherTok{=}  \FunctionTok{inla.spde2.pcmatern}\NormalTok{(mesh,}
                                  \AttributeTok{prior.sigma =} \FunctionTok{c}\NormalTok{(}\DecValTok{1}\NormalTok{, }\FloatTok{0.5}\NormalTok{),}
                                  \AttributeTok{prior.range =} \FunctionTok{c}\NormalTok{(}\DecValTok{100}\NormalTok{, }\FloatTok{0.5}\NormalTok{))}
\end{Highlighting}
\end{Shaded}

And here we plot the different priors for the range:

\begin{Shaded}
\begin{Highlighting}[]
\FunctionTok{ggplot}\NormalTok{() }\SpecialCharTok{+} 
  \FunctionTok{geom\_line}\NormalTok{(}\AttributeTok{data =} \FunctionTok{dens\_prior\_range}\NormalTok{(}\DecValTok{30}\NormalTok{,.}\DecValTok{5}\NormalTok{), }\FunctionTok{aes}\NormalTok{(x,y, }\AttributeTok{color =} \StringTok{"model1"}\NormalTok{)) }\SpecialCharTok{+}
  \FunctionTok{geom\_line}\NormalTok{(}\AttributeTok{data =} \FunctionTok{dens\_prior\_range}\NormalTok{(}\DecValTok{1000}\NormalTok{,.}\DecValTok{5}\NormalTok{), }\FunctionTok{aes}\NormalTok{(x,y, }\AttributeTok{color =} \StringTok{"model2"}\NormalTok{)) }\SpecialCharTok{+}
  \FunctionTok{geom\_line}\NormalTok{(}\AttributeTok{data =} \FunctionTok{dens\_prior\_range}\NormalTok{(}\DecValTok{100}\NormalTok{,.}\DecValTok{5}\NormalTok{), }\FunctionTok{aes}\NormalTok{(x,y, }\AttributeTok{color =} \StringTok{"model3"}\NormalTok{)) }
\end{Highlighting}
\end{Shaded}

\pandocbounded{\includegraphics[keepaspectratio]{Geostat_ex_files/figure-pdf/unnamed-chunk-12-1.pdf}}

and for the sd:

\begin{Shaded}
\begin{Highlighting}[]
\FunctionTok{ggplot}\NormalTok{() }\SpecialCharTok{+} 
  \FunctionTok{geom\_line}\NormalTok{(}\AttributeTok{data =} \FunctionTok{dens\_prior\_sd}\NormalTok{(}\DecValTok{1}\NormalTok{,.}\DecValTok{5}\NormalTok{), }\FunctionTok{aes}\NormalTok{(x,y, }\AttributeTok{color =} \StringTok{"model1"}\NormalTok{)) }\SpecialCharTok{+}
  \FunctionTok{geom\_line}\NormalTok{(}\AttributeTok{data =} \FunctionTok{dens\_prior\_sd}\NormalTok{(}\DecValTok{10}\NormalTok{,.}\DecValTok{5}\NormalTok{), }\FunctionTok{aes}\NormalTok{(x,y, }\AttributeTok{color =} \StringTok{"model2"}\NormalTok{)) }\SpecialCharTok{+}
  \FunctionTok{geom\_line}\NormalTok{(}\AttributeTok{data =} \FunctionTok{dens\_prior\_sd}\NormalTok{(.}\DecValTok{1}\NormalTok{,.}\DecValTok{5}\NormalTok{), }\FunctionTok{aes}\NormalTok{(x,y, }\AttributeTok{color =} \StringTok{"model3"}\NormalTok{)) }
\end{Highlighting}
\end{Shaded}

\pandocbounded{\includegraphics[keepaspectratio]{Geostat_ex_files/figure-pdf/unnamed-chunk-13-1.pdf}}

\begin{tcolorbox}[enhanced jigsaw, arc=.35mm, bottomrule=.15mm, bottomtitle=1mm, breakable, colback=white, colbacktitle=quarto-callout-tip-color!10!white, colframe=quarto-callout-tip-color-frame, coltitle=black, left=2mm, leftrule=.75mm, opacityback=0, opacitybacktitle=0.6, rightrule=.15mm, title={Question}, titlerule=0mm, toprule=.15mm, toptitle=1mm]

Consider the \texttt{pcod\_sf}, the spatial extension and type of the
data\ldots{} are some of the previous choices more reasonable than
others?

\begin{itemize}
\tightlist
\item
  \begin{enumerate}
  \def\labelenumi{(\Alph{enumi})}
  \tightlist
  \item
    spde\_model1\\
  \end{enumerate}
\item
  \begin{enumerate}
  \def\labelenumi{(\Alph{enumi})}
  \setcounter{enumi}{1}
  \tightlist
  \item
    spde\_model2\\
  \end{enumerate}
\item
  \begin{enumerate}
  \def\labelenumi{(\Alph{enumi})}
  \setcounter{enumi}{2}
  \tightlist
  \item
    spde\_model3
  \end{enumerate}
\end{itemize}

\textbf{NOTE} Remember that a prior should be reasonable, but the model
should not totally depend on it.

\end{tcolorbox}

\subsubsection{3. Define the components of the linear
predictor}\label{define-the-components-of-the-linear-predictor}

We have now defined a mesh and a SPDE representation of the spatial GF.
We now need to define the model components:

\begin{Shaded}
\begin{Highlighting}[]
\NormalTok{cmp }\OtherTok{=} \ErrorTok{\textasciitilde{}} \FunctionTok{Intercept}\NormalTok{(}\DecValTok{1}\NormalTok{) }\SpecialCharTok{+} \FunctionTok{space}\NormalTok{(geometry, }\AttributeTok{model =}\NormalTok{ spde\_model3)}
\end{Highlighting}
\end{Shaded}

\textbf{NOTE} since the dataframe we use (\texttt{pcod\_sf}) is an
\texttt{sf} object the input in the \texttt{space()} component is the
geometry of the dataset.

\subsubsection{4. Define the observation
model}\label{define-the-observation-model}

Our data are Bernoulli distributed so we can define the observation
model as:

\begin{Shaded}
\begin{Highlighting}[]
\NormalTok{formula }\OtherTok{=}\NormalTok{ present }\SpecialCharTok{\textasciitilde{}}\NormalTok{ Intercept  }\SpecialCharTok{+}\NormalTok{ space}

\NormalTok{lik }\OtherTok{=} \FunctionTok{bru\_obs}\NormalTok{(}\AttributeTok{formula =}\NormalTok{ formula, }
              \AttributeTok{data =}\NormalTok{ pcod\_sf, }
              \AttributeTok{family =} \StringTok{"binomial"}\NormalTok{)}
\end{Highlighting}
\end{Shaded}

\subsubsection{5. Run the model}\label{run-the-model}

Finally we are ready to run the model

\begin{Shaded}
\begin{Highlighting}[]
\NormalTok{fit1 }\OtherTok{=} \FunctionTok{bru}\NormalTok{(cmp,lik)}
\end{Highlighting}
\end{Shaded}

\subsection{Extract results}\label{extract-results}

\subsubsection{Hyperparameters}\label{hyperparameters}

\begin{tcolorbox}[enhanced jigsaw, arc=.35mm, bottomrule=.15mm, bottomtitle=1mm, breakable, colback=white, colbacktitle=quarto-callout-warning-color!10!white, colframe=quarto-callout-warning-color-frame, coltitle=black, left=2mm, leftrule=.75mm, opacityback=0, opacitybacktitle=0.6, rightrule=.15mm, title={Task}, titlerule=0mm, toprule=.15mm, toptitle=1mm]

Plot the posterior densities for the range \(\rho\) and the standard
deviation \(\sigma\) alog with the prior for both parameters.

Take hint

Posterior marginals for the hyperparameters can be extracted from the
fitted model as:

Click here to see the solution

\begin{Shaded}
\begin{Highlighting}[]
\NormalTok{fit1}\SpecialCharTok{$}\NormalTok{marginals.hyperpar}\SpecialCharTok{$}\StringTok{\textquotesingle{}name of the parameter\textquotesingle{}}
\end{Highlighting}
\end{Shaded}

Click here to see the solution

\begin{Shaded}
\begin{Highlighting}[]
\CommentTok{\# Extract marginal for the range}

\FunctionTok{ggplot}\NormalTok{() }\SpecialCharTok{+} 
  \FunctionTok{geom\_line}\NormalTok{(}\AttributeTok{data =}\NormalTok{ fit1}\SpecialCharTok{$}\NormalTok{marginals.hyperpar}\SpecialCharTok{$}\StringTok{\textasciigrave{}}\AttributeTok{Range for space}\StringTok{\textasciigrave{}}\NormalTok{,}
            \FunctionTok{aes}\NormalTok{(x,y, }\AttributeTok{color =} \StringTok{"Posterior"}\NormalTok{)) }\SpecialCharTok{+}
  \FunctionTok{geom\_line}\NormalTok{(}\AttributeTok{data =} \FunctionTok{dens\_prior\_range}\NormalTok{(}\DecValTok{100}\NormalTok{,.}\DecValTok{5}\NormalTok{),}
            \FunctionTok{aes}\NormalTok{(x,y, }\AttributeTok{color =} \StringTok{"Prior"}\NormalTok{))}


\FunctionTok{ggplot}\NormalTok{() }\SpecialCharTok{+} 
  \FunctionTok{geom\_line}\NormalTok{(}\AttributeTok{data =}\NormalTok{ fit1}\SpecialCharTok{$}\NormalTok{marginals.hyperpar}\SpecialCharTok{$}\StringTok{\textasciigrave{}}\AttributeTok{Stdev for space}\StringTok{\textasciigrave{}}\NormalTok{,}
            \FunctionTok{aes}\NormalTok{(x,y, }\AttributeTok{color =} \StringTok{"Posterior"}\NormalTok{)) }\SpecialCharTok{+}
  \FunctionTok{geom\_line}\NormalTok{(}\AttributeTok{data =} \FunctionTok{dens\_prior\_sd}\NormalTok{(}\DecValTok{1}\NormalTok{,.}\DecValTok{5}\NormalTok{), }\FunctionTok{aes}\NormalTok{(x,y, }\AttributeTok{color =} \StringTok{"Prior"}\NormalTok{))}
\end{Highlighting}
\end{Shaded}

\end{tcolorbox}

\subsection{Spatial prediction}\label{spatial-prediction}

We now want to extract the estimated posterior mean and sd of spatial
GF. To do this we first need to define a grid of points where we want to
predict. We do this using the function \texttt{fm\_pixel()}, which
creates a regular grid of points covering the mesh

\begin{Shaded}
\begin{Highlighting}[]
\NormalTok{pxl }\OtherTok{=} \FunctionTok{fm\_pixels}\NormalTok{(mesh)}
\end{Highlighting}
\end{Shaded}

then compute the prediction for both the spatial GF and the linear
predictor (spatial GF + intercept)

\begin{Shaded}
\begin{Highlighting}[]
\NormalTok{preds }\OtherTok{=} \FunctionTok{predict}\NormalTok{(fit1, pxl, }\SpecialCharTok{\textasciitilde{}}\FunctionTok{data.frame}\NormalTok{(}\AttributeTok{spatial =}\NormalTok{ space,}
                                      \AttributeTok{total =}\NormalTok{ Intercept }\SpecialCharTok{+}\NormalTok{ space))}
\end{Highlighting}
\end{Shaded}

Finally, we can plot the maps

\begin{Shaded}
\begin{Highlighting}[]
\FunctionTok{ggplot}\NormalTok{() }\SpecialCharTok{+} \FunctionTok{geom\_sf}\NormalTok{(}\AttributeTok{data =}\NormalTok{ preds}\SpecialCharTok{$}\NormalTok{spatial,}\FunctionTok{aes}\NormalTok{(}\AttributeTok{color =}\NormalTok{ mean)) }\SpecialCharTok{+} \FunctionTok{scale\_color\_scico}\NormalTok{() }\SpecialCharTok{+} \FunctionTok{ggtitle}\NormalTok{(}\StringTok{"Posterior mean"}\NormalTok{)}
\end{Highlighting}
\end{Shaded}

\pandocbounded{\includegraphics[keepaspectratio]{Geostat_ex_files/figure-pdf/unnamed-chunk-21-1.pdf}}

\begin{Shaded}
\begin{Highlighting}[]
\FunctionTok{ggplot}\NormalTok{() }\SpecialCharTok{+} \FunctionTok{geom\_sf}\NormalTok{(}\AttributeTok{data =}\NormalTok{ preds}\SpecialCharTok{$}\NormalTok{spatial,}\FunctionTok{aes}\NormalTok{(}\AttributeTok{color =}\NormalTok{ sd)) }\SpecialCharTok{+} \FunctionTok{scale\_color\_scico}\NormalTok{() }\SpecialCharTok{+} \FunctionTok{ggtitle}\NormalTok{(}\StringTok{"Posterior sd"}\NormalTok{)}
\end{Highlighting}
\end{Shaded}

\pandocbounded{\includegraphics[keepaspectratio]{Geostat_ex_files/figure-pdf/unnamed-chunk-21-2.pdf}}

\textbf{Note} The posterior sd is lowest at the observation points. Note
how the posterior sd is inflated around the border, this is the ``border
effect'' due to the SPDE representation.

Instead of predicting over a grid covering the whole mesh, we can limit
our predictions to the points where the covariate is defined. We can do
this by defining a \texttt{sf} object using coordinates in the object
\texttt{depth\_r}.

\begin{Shaded}
\begin{Highlighting}[]
\NormalTok{pxl1 }\OtherTok{=} \FunctionTok{data.frame}\NormalTok{(}\FunctionTok{crds}\NormalTok{(depth\_r), }
                  \FunctionTok{as.data.frame}\NormalTok{(depth\_r}\SpecialCharTok{$}\NormalTok{depth)) }\SpecialCharTok{\%\textgreater{}\%} 
       \FunctionTok{filter}\NormalTok{(}\SpecialCharTok{!}\FunctionTok{is.na}\NormalTok{(depth)) }\SpecialCharTok{\%\textgreater{}\%}
\FunctionTok{st\_as\_sf}\NormalTok{(}\AttributeTok{coords =} \FunctionTok{c}\NormalTok{(}\StringTok{"x"}\NormalTok{,}\StringTok{"y"}\NormalTok{))}
\end{Highlighting}
\end{Shaded}

\begin{tcolorbox}[enhanced jigsaw, arc=.35mm, bottomrule=.15mm, bottomtitle=1mm, breakable, colback=white, colbacktitle=quarto-callout-warning-color!10!white, colframe=quarto-callout-warning-color-frame, coltitle=black, left=2mm, leftrule=.75mm, opacityback=0, opacitybacktitle=0.6, rightrule=.15mm, title={Task}, titlerule=0mm, toprule=.15mm, toptitle=1mm]

Produce prediction over \texttt{pxl1} using the same techniques as
before. Plot your results.

Take hint

Add hint details here\ldots{}

Click here to see the solution

\begin{Shaded}
\begin{Highlighting}[]
\NormalTok{pred\_pxl1 }\OtherTok{=} \FunctionTok{predict}\NormalTok{(fit1, pxl1, }\SpecialCharTok{\textasciitilde{}}\FunctionTok{data.frame}\NormalTok{(}\AttributeTok{spatial =}\NormalTok{ space,}
                                      \AttributeTok{total =}\NormalTok{ Intercept }\SpecialCharTok{+}\NormalTok{ space))}

\FunctionTok{ggplot}\NormalTok{() }\SpecialCharTok{+} \FunctionTok{geom\_sf}\NormalTok{(}\AttributeTok{data =}\NormalTok{ pred\_pxl1}\SpecialCharTok{$}\NormalTok{total,}\FunctionTok{aes}\NormalTok{(}\AttributeTok{color =}\NormalTok{ mean)) }\SpecialCharTok{+} \FunctionTok{scale\_color\_scico}\NormalTok{() }\SpecialCharTok{+} \FunctionTok{ggtitle}\NormalTok{(}\StringTok{"Posterior mean"}\NormalTok{)}
\end{Highlighting}
\end{Shaded}

\pandocbounded{\includegraphics[keepaspectratio]{Geostat_ex_files/figure-pdf/unnamed-chunk-23-1.pdf}}

\begin{Shaded}
\begin{Highlighting}[]
\FunctionTok{ggplot}\NormalTok{() }\SpecialCharTok{+} \FunctionTok{geom\_sf}\NormalTok{(}\AttributeTok{data =}\NormalTok{ pred\_pxl1}\SpecialCharTok{$}\NormalTok{total,}\FunctionTok{aes}\NormalTok{(}\AttributeTok{color =}\NormalTok{ sd)) }\SpecialCharTok{+} \FunctionTok{scale\_color\_scico}\NormalTok{() }\SpecialCharTok{+} \FunctionTok{ggtitle}\NormalTok{(}\StringTok{"Posterior sd"}\NormalTok{)}
\end{Highlighting}
\end{Shaded}

\pandocbounded{\includegraphics[keepaspectratio]{Geostat_ex_files/figure-pdf/unnamed-chunk-23-2.pdf}}

\end{tcolorbox}

Instead of computing the posterior mean and standard deviations of the
estimated surface, one can also \emph{simulate} possible realizations of
such a surface. This will give the user a better idea of the type of
realized surfaces one can expect. We can do this using the function
\texttt{generate()}.

\begin{Shaded}
\begin{Highlighting}[]
\CommentTok{\# we simulate 4 samples from the }
\NormalTok{gens }\OtherTok{=} \FunctionTok{generate}\NormalTok{(fit1, pxl1, }\SpecialCharTok{\textasciitilde{}}\NormalTok{ (Intercept }\SpecialCharTok{+}\NormalTok{ space),}
                \AttributeTok{n.samples =} \DecValTok{4}\NormalTok{)}

\NormalTok{pp }\OtherTok{=} \FunctionTok{cbind}\NormalTok{(pxl1, gens)}

\NormalTok{pp }\SpecialCharTok{\%\textgreater{}\%} \FunctionTok{select}\NormalTok{(}\SpecialCharTok{{-}}\NormalTok{depth) }\SpecialCharTok{\%\textgreater{}\%}
  \FunctionTok{pivot\_longer}\NormalTok{(}\SpecialCharTok{{-}}\NormalTok{geometry) }\SpecialCharTok{\%\textgreater{}\%}
    \FunctionTok{ggplot}\NormalTok{() }\SpecialCharTok{+} 
      \FunctionTok{geom\_sf}\NormalTok{(}\FunctionTok{aes}\NormalTok{(}\AttributeTok{color =}\NormalTok{ value)) }\SpecialCharTok{+}
      \FunctionTok{facet\_wrap}\NormalTok{(.}\SpecialCharTok{\textasciitilde{}}\NormalTok{name) }\SpecialCharTok{+}
        \FunctionTok{scale\_color\_scico}\NormalTok{(}\AttributeTok{direction =} \SpecialCharTok{{-}}\DecValTok{1}\NormalTok{) }\SpecialCharTok{+}
        \FunctionTok{ggtitle}\NormalTok{(}\StringTok{"Sample from the fitted model"}\NormalTok{)}
\end{Highlighting}
\end{Shaded}

\pandocbounded{\includegraphics[keepaspectratio]{Geostat_ex_files/figure-pdf/unnamed-chunk-24-1.pdf}}

\subsection{An alternative model}\label{an-alternative-model}

We now want to check if the \texttt{depth} covatiate has an influence on
the probability of presence. We do this in two different models

\begin{enumerate}
\def\labelenumi{\arabic{enumi}.}
\tightlist
\item
  \textbf{Model 1} The depth enters the model in a linear way. The
  linear predictor is then defined as:
\end{enumerate}

\[
  \eta(s) = \text{logit}(p(s)) = \beta_0 + \omega(s) + \beta_1\ \text{depth}(s)
\]

\begin{enumerate}
\def\labelenumi{\arabic{enumi}.}
\setcounter{enumi}{1}
\tightlist
\item
  \textbf{Model 1} The depth enters the model in a non linear way. The
  linear predictor is then defined as:
\end{enumerate}

\[
  \eta(s) = \text{logit}(p(s)) = \beta_0 + \omega(s) +  f(\text{depth}(s))
\] where \(f(.)\) is a smooth function. We will use a RW2 model for
this.

\begin{tcolorbox}[enhanced jigsaw, arc=.35mm, bottomrule=.15mm, bottomtitle=1mm, breakable, colback=white, colbacktitle=quarto-callout-warning-color!10!white, colframe=quarto-callout-warning-color-frame, coltitle=black, left=2mm, leftrule=.75mm, opacityback=0, opacitybacktitle=0.6, rightrule=.15mm, title={Task}, titlerule=0mm, toprule=.15mm, toptitle=1mm]

Fit model 1. Define components, observation model and use the
\texttt{bru()} function to estimate the parameters.

\textbf{Note} Use the scaled version of the covariate stored in
\texttt{depth\_r\$depth\_scaled}.

What is the liner effect of depth on the logit probability?

Take hint

Add hint details here\ldots{}

Click here to see the solution

\begin{Shaded}
\begin{Highlighting}[]
\NormalTok{cmp }\OtherTok{=} \ErrorTok{\textasciitilde{}} \FunctionTok{Intercept}\NormalTok{(}\DecValTok{1}\NormalTok{) }\SpecialCharTok{+} \FunctionTok{space}\NormalTok{(geometry, }\AttributeTok{model =}\NormalTok{ spde\_model3) }\SpecialCharTok{+}
        \FunctionTok{covariate}\NormalTok{(depth\_r}\SpecialCharTok{$}\NormalTok{depth\_scaled, }\AttributeTok{model =} \StringTok{"linear"}\NormalTok{)}

\NormalTok{formula }\OtherTok{=}\NormalTok{ present }\SpecialCharTok{\textasciitilde{}}\NormalTok{ Intercept  }\SpecialCharTok{+}\NormalTok{ space }\SpecialCharTok{+}\NormalTok{ covariate}

\NormalTok{lik }\OtherTok{=} \FunctionTok{bru\_obs}\NormalTok{(}\AttributeTok{formula =}\NormalTok{ formula, }
              \AttributeTok{data =}\NormalTok{ pcod\_sf, }
              \AttributeTok{family =} \StringTok{"binomial"}\NormalTok{)}


\NormalTok{fit2 }\OtherTok{=} \FunctionTok{bru}\NormalTok{(cmp, lik)}
\end{Highlighting}
\end{Shaded}

\end{tcolorbox}

We now want to fit \textbf{Model 2} where we allow the effect of depth
to be non-linear. To use the RW2 model we need to \emph{group} the
values of depth into distinct classe. To do this we use the function
\texttt{inla.group()} which, by default, creates 20 groups. The we can
fit the model as usual

\begin{Shaded}
\begin{Highlighting}[]
\CommentTok{\# create the grouped variable}
\NormalTok{depth\_r}\SpecialCharTok{$}\NormalTok{depth\_group }\OtherTok{=} \FunctionTok{inla.group}\NormalTok{(}\FunctionTok{values}\NormalTok{(depth\_r}\SpecialCharTok{$}\NormalTok{depth\_scaled))}

\CommentTok{\# run the model}
\NormalTok{cmp }\OtherTok{=} \ErrorTok{\textasciitilde{}} \FunctionTok{Intercept}\NormalTok{(}\DecValTok{1}\NormalTok{) }\SpecialCharTok{+} \FunctionTok{space}\NormalTok{(geometry, }\AttributeTok{model =}\NormalTok{ spde\_model3) }\SpecialCharTok{+}
        \FunctionTok{covariate}\NormalTok{(depth\_r}\SpecialCharTok{$}\NormalTok{depth\_group, }\AttributeTok{model =} \StringTok{"rw2"}\NormalTok{)}

\NormalTok{formula }\OtherTok{=}\NormalTok{ present }\SpecialCharTok{\textasciitilde{}}\NormalTok{ Intercept  }\SpecialCharTok{+}\NormalTok{ space }\SpecialCharTok{+}\NormalTok{ covariate}

\NormalTok{lik }\OtherTok{=} \FunctionTok{bru\_obs}\NormalTok{(}\AttributeTok{formula =}\NormalTok{ formula, }
              \AttributeTok{data =}\NormalTok{ pcod\_sf, }
              \AttributeTok{family =} \StringTok{"binomial"}\NormalTok{)}


\NormalTok{fit3 }\OtherTok{=} \FunctionTok{bru}\NormalTok{(cmp, lik)}

\CommentTok{\# plot the estimated effect of depth}

\NormalTok{fit3}\SpecialCharTok{$}\NormalTok{summary.random}\SpecialCharTok{$}\NormalTok{covariate }\SpecialCharTok{\%\textgreater{}\%} 
  \FunctionTok{ggplot}\NormalTok{() }\SpecialCharTok{+} \FunctionTok{geom\_line}\NormalTok{(}\FunctionTok{aes}\NormalTok{(ID,mean)) }\SpecialCharTok{+} 
                                  \FunctionTok{geom\_ribbon}\NormalTok{(}\FunctionTok{aes}\NormalTok{(ID, }\AttributeTok{ymin =} \StringTok{\textasciigrave{}}\AttributeTok{0.025quant}\StringTok{\textasciigrave{}}\NormalTok{, }
                                                      \AttributeTok{ymax =} \StringTok{\textasciigrave{}}\AttributeTok{0.975quant}\StringTok{\textasciigrave{}}\NormalTok{), }\AttributeTok{alpha =} \FloatTok{0.5}\NormalTok{)}
\end{Highlighting}
\end{Shaded}

\pandocbounded{\includegraphics[keepaspectratio]{Geostat_ex_files/figure-pdf/unnamed-chunk-26-1.pdf}}

\begin{tcolorbox}[enhanced jigsaw, arc=.35mm, bottomrule=.15mm, bottomtitle=1mm, breakable, colback=white, colbacktitle=quarto-callout-warning-color!10!white, colframe=quarto-callout-warning-color-frame, coltitle=black, left=2mm, leftrule=.75mm, opacityback=0, opacitybacktitle=0.6, rightrule=.15mm, title={Task}, titlerule=0mm, toprule=.15mm, toptitle=1mm]

Create a map of predicted \emph{probability} from Model 3. You can use a
inverse logit function defined as

\begin{Shaded}
\begin{Highlighting}[]
\NormalTok{inv\_logit }\OtherTok{=} \ControlFlowTok{function}\NormalTok{(x) (}\DecValTok{1}\SpecialCharTok{+}\FunctionTok{exp}\NormalTok{(}\SpecialCharTok{{-}}\NormalTok{x))}\SpecialCharTok{\^{}}\NormalTok{(}\SpecialCharTok{{-}}\DecValTok{1}\NormalTok{)}
\end{Highlighting}
\end{Shaded}

Take hint

The \texttt{predict()} function can take as input also functions of
elements of the components you want to consider

Click here to see the solution

\begin{Shaded}
\begin{Highlighting}[]
\NormalTok{inv\_logit }\OtherTok{=} \ControlFlowTok{function}\NormalTok{(x) (}\DecValTok{1}\SpecialCharTok{+}\FunctionTok{exp}\NormalTok{(}\SpecialCharTok{{-}}\NormalTok{x))}\SpecialCharTok{\^{}}\NormalTok{(}\SpecialCharTok{{-}}\DecValTok{1}\NormalTok{)}

\NormalTok{pred3  }\OtherTok{=} \FunctionTok{predict}\NormalTok{(fit3, pxl1, }\SpecialCharTok{\textasciitilde{}}\FunctionTok{inv\_logit}\NormalTok{(Intercept }\SpecialCharTok{+}\NormalTok{ space }\SpecialCharTok{+}\NormalTok{ covariate) )}

\NormalTok{pred3 }\SpecialCharTok{\%\textgreater{}\%} \FunctionTok{ggplot}\NormalTok{() }\SpecialCharTok{+} 
      \FunctionTok{geom\_sf}\NormalTok{(}\FunctionTok{aes}\NormalTok{(}\AttributeTok{color =}\NormalTok{ mean)) }\SpecialCharTok{+}
        \FunctionTok{scale\_color\_scico}\NormalTok{(}\AttributeTok{direction =} \SpecialCharTok{{-}}\DecValTok{1}\NormalTok{) }\SpecialCharTok{+}
        \FunctionTok{ggtitle}\NormalTok{(}\StringTok{"Sample from the fitted model"}\NormalTok{)}
\end{Highlighting}
\end{Shaded}

\pandocbounded{\includegraphics[keepaspectratio]{Geostat_ex_files/figure-pdf/unnamed-chunk-28-1.pdf}}

\end{tcolorbox}




\end{document}
